\section*{Capítulo \ref{ch:b2:liquid-vapour-diagrams} --- Diagramas binarios líquido--vapor y destilación}

\begin{solution}{exo:b2:liquid-vapour-diagrams:1}
Benceno \qty{80.1}{\degreeCelsius}, tolueno \qty{110.6}{\degreeCelsius}. El líquido
de composición 0.2 empieza a hervir hacia \qty{102}{\degreeCelsius}; su primer
vapor tiene una composición de 0.38.
\end{solution}

\begin{solution}{exo:b2:liquid-vapour-diagrams:2}
$p = 0.5 \times 1.010 + 0.5 \times 0.388 = \qty{0.699}{bar}$;
$y = 0.505/0.699 = 0.722$.
\end{solution}

\begin{solution}{exo:b2:liquid-vapour-diagrams:3}
$n_V/n = (0.30 - 0.20)/(0.45 - 0.20) = 0.40$: \qty{4.0}{mol} de vapor y
\qty{6.0}{mol} de líquido. Comprobación: $6.0 \times 0.20 + 4.0 \times 0.45 = 3.0 = 10
\times 0.30$.
\end{solution}

\begin{solution}{exo:b2:liquid-vapour-diagrams:4}
(a) $v = 2 + 2 - 1 = 3$: $T$, $p$ y $x$ son libres. (b) $v = 2 + 2 - 2 = 2$; a
presión fija, 1. (c) El \omterm{def:b2:liquid-vapour-diagrams:azeotrope}{azeótropo} añade la relación $x = y$: $v = 1$; a presión
fija, 0: el \omterm{def:b2:liquid-vapour-diagrams:azeotrope}{azeótropo} hierve a temperatura fija, como un líquido puro,
aunque su composición cambia con la presión.
\end{solution}

\begin{solution}{exo:b2:liquid-vapour-diagrams:5}
En el \omterm{def:b2:liquid-vapour-diagrams:bubble-dew}{punto de rocío}, $yp/p_1^* + (1 - y)p/p_2^* = 1$. A \qty{100}{\degreeCelsius}:
$0.3 \times 1.013/1.80 + 0.7 \times 1.013/0.742 = 1.124 > 1$ (demasiado frío); a
\qty{105}{\degreeCelsius}: $0.148 + 0.824 = 0.971 < 1$. Interpolando, $T \approx
100 + 5 \times 0.124/0.153 = \qty{104}{\degreeCelsius}$. La primera gota tiene $x =
yp/p_1^* \approx 0.3 \times 1.013/2.0 = 0.15$.
\end{solution}

\begin{solution}{exo:b2:liquid-vapour-diagrams:6}
Las primeras gotas tienen la composición del vapor que hay sobre el líquido en ebullición,
$y = 0.71$, a \qty{92.1}{\degreeCelsius}. Como el benceno sale preferentemente, el
líquido se empobrece en él, su punto de ebullición sube y el destilado también se
empobrece; las últimas gotas son tolueno casi puro. La destilación simple es una
sola etapa de equilibrio, repetida sobre un líquido que no deja de cambiar: cada
fracción es una mezcla.
\end{solution}

\begin{solution}{exo:b2:liquid-vapour-diagrams:7}
Alimentación de 0.05, del lado del agua respecto del \omterm{def:b2:liquid-vapour-diagrams:azeotrope}{azeótropo}: la cabeza proporciona el \omterm{def:b2:liquid-vapour-diagrams:azeotrope}{azeótropo}
(0.88 en moles, 95~\% en masa) y el hervidor retiene el agua. Alimentación de 0.95, del lado
del etanol: la cabeza sigue proporcionando el \omterm{def:b2:liquid-vapour-diagrams:azeotrope}{azeótropo}, el punto de ebullición más bajo,
y el hervidor retiene etanol puro.
\end{solution}

\begin{solution}{exo:b2:liquid-vapour-diagrams:8}
La mezcla hierve cuando $0.10 + p^*_{\text{agua}} = 1.013$, es decir,
$p^*_{\text{agua}} = \qty{0.91}{bar}$, hacia \qty{97}{\degreeCelsius}. Razón de masas
$0.10 \times 136/(0.91 \times 18.02) = 0.83$: \qty{0.83}{g} de esencia por gramo
de agua.
\end{solution}

\begin{solution}{exo:b2:liquid-vapour-diagrams:9}
A temperatura ambiente, dos capas L$_2$ y L$_1$ (la composición 0.4 está en la
laguna). Al calentar, las capas hierven juntas a la temperatura del \omterm{def:b2:liquid-vapour-diagrams:miscibility-gap}{heteroazeótropo},
que se mantiene constante; el primer vapor tiene la composición
heteroazeotrópica (el punto). Como 0.4 está del lado de L$_2$ respecto del punto, la capa
L$_1$ se agota primero; la temperatura sube entonces, la L$_2$ restante se
empobrece en 1 a lo largo de su \omterm{def:b2:liquid-vapour-diagrams:bubble-dew}{curva de burbuja} y el vapor sigue la \omterm{def:b2:liquid-vapour-diagrams:bubble-dew}{curva de rocío}, hasta
que el punto alcanza la \omterm{def:b2:liquid-vapour-diagrams:bubble-dew}{curva de rocío}: se evapora la última gota.
\end{solution}

\begin{solution}{exo:b2:liquid-vapour-diagrams:10}
$N_{\min} = \ln(99 \times 99)/\ln 2.47 = 9.19/0.904 = 10.2$, frente a 6.5: pasar
de un 95~\% a un 99~\% de pureza en ambos extremos cuesta más de la mitad de platos
adicionales. Cada factor adicional en $x/(1 - x)$ cuesta el mismo número de platos, y
la pureza se mide por la impureza que queda, que solo disminuye geométricamente.
\end{solution}

\begin{solution}{exo:b2:liquid-vapour-diagrams:11}
El \omterm{def:b2:liquid-vapour-diagrams:binary-diagram}{diagrama isobárico} solo no basta: hacen falta $p_1^*$ y $p_2^*$ a
\qty{80}{\degreeCelsius}, a partir de las ecuaciones de Antoine (1.010 y
\qty{0.388}{bar}). Entonces $p = p_2^* + x(p_1^* - p_2^*)$ y $p = 1/(y/p_1^* + (1 -
y)/p_2^*)$. El líquido es estable a presión alta (comprimir un vapor
lo condensa), y a temperatura baja en el \omterm{def:b2:liquid-vapour-diagrams:binary-diagram}{diagrama isobárico}: los dos diagramas
están invertidos el uno respecto del otro.
\end{solution}

\begin{solution}{exo:b2:liquid-vapour-diagrams:12}
Como $y$ crece con $x$ e $y(1 - y) > 0$, $dp/dx$ tiene el signo de $y - x$:
la presión aumenta con $x$ donde el vapor es más rico en 1. Para una \omterm{def:b2:chemical-potential:ideal-mixture}{mezcla
ideal}, $y > x$ para todo $0 < x < 1$ (\cref{prop:b2:liquid-vapour-diagrams:richer}):
$p$ es estrictamente monótona, no tiene extremo y no hay \omterm{def:b2:liquid-vapour-diagrams:azeotrope}{azeótropo}.
\end{solution}

\begin{solution}{pb:b2:liquid-vapour-diagrams:1}
\textbf{1.} $T = B/(A - \log_{10}1.013) - C$: benceno $1203.835/4.01253 + 53.226 =
\qty{353.2}{K}$, \qty{80.1}{\degreeCelsius}; tolueno $1343.943/4.07266 + 53.773 =
\qty{383.8}{K}$, \qty{110.6}{\degreeCelsius}.
\textbf{2.} A \qty{363.15}{K}: $p_1^* = \qty{1.361}{bar}$, $p_2^* = \qty{0.542}{bar}$.
\textbf{3.} $x = (1.013 - 0.542)/(1.361 - 0.542) = 0.575$; $y = 0.575 \times
1.361/1.013 = 0.773$.
\textbf{4.} $x = (1.013 - 0.742)/(1.800 - 0.742) = 0.256$; $y = 0.256 \times
1.800/1.013 = 0.455$.
\textbf{5.} \omterm{def:b2:liquid-vapour-diagrams:bubble-dew}{Puntos de burbuja} (0, 110.6), (0.256, 100), (0.575, 90), (1, 80.1); \omterm{def:b2:liquid-vapour-diagrams:bubble-dew}{puntos de
rocío} (0, 110.6), (0.455, 100), (0.773, 90), (1, 80.1): la lente de la
figura.
\textbf{6.} $\frac12(p_1^* + p_2^*) = 1.013$: a \qty{90}{\degreeCelsius} 0.952, a
\qty{95}{\degreeCelsius} 1.102; interpolando, $90 + 5 \times 0.061/0.150 \approx
\qty{92}{\degreeCelsius}$ (\qty{92.1}{\degreeCelsius} exactamente).
\textbf{7.} $v = 2$, 1 a presión fija: cada temperatura fija $x$ e $y$,
y por eso el diagrama está formado por dos curvas.
\textbf{8.} $x = (1.013 - 0.636)/(1.569 - 0.636) = 0.404$; $y = 0.404 \times
1.569/1.013 = 0.626$.
\textbf{9.} $n_V = 100 \times (0.500 - 0.404)/(0.626 - 0.404) = \qty{43}{mol}$,
$n_L = \qty{57}{mol}$.
\textbf{10.} Vapor $43 \times 0.626 = \qty{26.9}{mol}$; líquido $57 \times 0.404
= \qty{23.0}{mol}$; total \qty{49.9}{mol} $\approx 50$.
\textbf{11.} $26.9/50 = 54~\%$ del benceno, en un vapor de solo un 63~\%.
\textbf{12.} A \qty{100}{\degreeCelsius} el vapor en equilibrio tiene $y =
0.455 < 0.5$: el punto de la alimentación está por encima de la \omterm{def:b2:liquid-vapour-diagrams:bubble-dew}{curva de rocío}, la alimentación es toda
vapor y no se separa nada.
\textbf{13.} En su \omterm{def:b2:liquid-vapour-diagrams:bubble-dew}{punto de rocío}, \qty{98.8}{\degreeCelsius} (por el método del
\cref{exo:b2:liquid-vapour-diagrams:5}).
\textbf{14.} $\alpha = p_1^*/p_2^*$; a \qty{80.1}{\degreeCelsius}, $1.013/0.390 =
2.60$.
\textbf{15.} A \qty{110.6}{\degreeCelsius}, $2.377/1.013 = 2.35$.
\textbf{16.} $\sqrt{2.60 \times 2.35} = 2.47$.
\textbf{17.} $y = xp_1^*/p$ y $1 - y = (1 - x)p_2^*/p$; se dividen.
\textbf{18.} No se extrae nada: entre dos platos el vapor que sube y
el líquido que baja transportan la misma cantidad de materia y de cada componente,
$V y_k = L x_{k+1}$ con $V = L$, luego $x_{k+1} = y_k$.
\textbf{19.} Con $r = x/(1 - x)$: $r(y_k) = \alpha\,r(x_k)$ y $x_{k+1} = y_k$, luego
$r(x_{k+1}) = \alpha\,r(x_k)$ y $r(x_D) = \alpha^N r(x_B)$; $N = \ln[r(x_D)/
r(x_B)]/\ln\alpha$.
\textbf{20.} $N_{\min} = \ln(19 \times 19)/\ln 2.47 = 5.889/0.904 = 6.5$.
\textbf{21.} Una columna debe suministrar producto, así que funciona con una \omterm{def:b2:liquid-vapour-diagrams:fractional-distillation}{razón de reflujo} finita,
que exige más platos; y un plato real no alcanza el equilibrio (su
eficacia es inferior a uno).
\textbf{22.} Siete peldaños, el último terminando en 0.034, por debajo de 0.05: 6.5 platos en un
recuento continuo, siete enteros (contando el hervidor como uno).
\textbf{23.} $x = (95/46.07)/(95/46.07 + 5/18.02) = 0.881$.
\textbf{24.} En la cabeza, como mucho el \omterm{def:b2:liquid-vapour-diagrams:azeotrope}{azeótropo} (0.88 en moles); en el fondo,
agua. Cincuenta platos no atraviesan el \omterm{def:b2:liquid-vapour-diagrams:azeotrope}{azeótropo}.
\textbf{25.} A reflujo total, $\boldsymbol{N_{\min} \approx 6.5}$ \omterm{def:b2:liquid-vapour-diagrams:fractional-distillation}{platos
teóricos}.
\end{solution}
